\section*{الفصل \ref{ch:g10:the-mole} --- المول والكتلة المولية}

\begin{solution}{exo:g10:the-mole:1}
$M(\ce{O2}) = 2 \times 16.0 = \qty{32.0}{g/mol}$;
$M(\ce{N2}) = 2 \times 14.0 = \qty{28.0}{g/mol}$;
$M(\ce{CH4}) = 12.0 + 4 \times 1.0 = \qty{16.0}{g/mol}$;
$M(\ce{NH3}) = 14.0 + 3 \times 1.0 = \qty{17.0}{g/mol}$;
$M(\ce{CO2}) = 12.0 + 2 \times 16.0 = \qty{44.0}{g/mol}$.
\end{solution}

\begin{solution}{exo:g10:the-mole:2}
الماء: $n = 36.0 / 18.0 = \qty{2.00}{mol}$. الميثان:
$n = 4.0 / 16.0 = \qty{0.25}{mol}$.
\end{solution}

\begin{solution}{exo:g10:the-mole:3}
$m(\ce{NaCl}) = 0.25 \times 58.5 = \qty{14.6}{g}$، أي نحو \qty{15}{g}؛
و$m(\ce{CO2}) = 2.0 \times 44.0 = \qty{88}{g}$.
\end{solution}

\begin{solution}{exo:g10:the-mole:4}
$N = 0.50 \times \num{6.02e23} = \num{3.0e23}$ \omterm{def:g7:atoms-and-molecules:molecule}{جزيء} من ثنائي الأكسجين.
و$n = \num{3.0e24} / \qty{6.02e23}{mol^{-1}} = \qty{5.0}{mol}$ من الحديد.
\end{solution}

\begin{solution}{exo:g10:the-mole:5}
$V = n \times V_m = 0.50 \times 24.1 = \qty{12}{L}$.
$n = V / V_m = 12 / 24.1 = \qty{0.50}{mol}$.
\end{solution}

\begin{solution}{exo:g10:the-mole:6}
كتلة القطرة \qty{0.050}{g}، إذن
$n = 0.050 / 18.0 = \qty{2.8e-3}{mol}$
و$N = \num{2.8e-3} \times \num{6.02e23} = \num{1.7e21}$ \omterm{def:g7:atoms-and-molecules:molecule}{جزيء}.
\end{solution}

\begin{solution}{exo:g10:the-mole:7}
\omterm{def:g7:atoms-and-molecules:atom}{ذرة} الألومنيوم (\qty{27.0}{g/mol}) نحو نصف \omterm{def:g7:atoms-and-molecules:atom}{ذرة} الحديد
(\qty{55.8}{g/mol}) ثقلًا، فالكتلة نفسها من الألومنيوم تحوي نحو ضعف عدد
\omterm{def:g7:atoms-and-molecules:atom}{الذرات}. التحقق: الألومنيوم $10/27.0 = \qty{0.370}{mol}$، أي
\num{2.23e23} \omterm{def:g7:atoms-and-molecules:atom}{ذرة}؛ والحديد $10/55.8 = \qty{0.179}{mol}$، أي
\num{1.08e23} \omterm{def:g7:atoms-and-molecules:atom}{ذرة}. فالألومنيوم أكثر، بعامل يقارب 2.1.
\end{solution}

\begin{solution}{exo:g10:the-mole:8}
اللتر الواحد $1/24.1 = \qty{0.0415}{mol}$ من الغاز. ثنائي أكسيد الكربون:
$0.0415 \times 44.0 = \qty{1.83}{g}$؛ ثنائي النيتروجين:
$0.0415 \times 28.0 = \qty{1.16}{g}$. فثنائي أكسيد الكربون أكثف بنحو مرة
ونصف من \omterm{def:g4:air-a-mixture-of-gases:air}{الهواء}، الذي تقارب كتلته المولية كتلة ثنائي
النيتروجين. والعصير المتخمّر يطلق ثنائي أكسيد الكربون، فيهبط ويتراكم قرب أرضية
القبو المغلق، حيث قد يخنق كل من يدخله.
\end{solution}

\begin{solution}{exo:g10:the-mole:9}
$M(\ce{C6H12O6}) = 6 \times 12.0 + 12 \times 1.0 + 6 \times 16.0 =
\qty{180.0}{g/mol}$. السهم $\div N_A$: $n = \num{1.0e22} / \num{6.02e23} =
\qty{1.66e-2}{mol}$؛ والسهم $\times M$: $m = \num{1.66e-2} \times 180.0 =
\qty{3.0}{g}$.
\end{solution}

\begin{solution}{exo:g10:the-mole:10}
$M = \qty{342.0}{g/mol}$، إذن $n = 5.0 / 342.0 = \qty{1.46e-2}{mol}$؛
و$N = \num{1.46e-2} \times \num{6.02e23} = \num{8.8e21}$ \omterm{def:g7:atoms-and-molecules:molecule}{جزيء}. وفي كل
منها 12 \omterm{def:g7:atoms-and-molecules:atom}{ذرة} كربون: $12 \times \num{8.8e21} = \num{1.1e23}$ \omterm{def:g7:atoms-and-molecules:atom}{ذرة} كربون.
\end{solution}

\begin{solution}{exo:g10:the-mole:11}
الماء: $1000 / 18.0 = \qty{55.6}{mol}$. الإيثانول:
$M = 2 \times 12.0 + 6 \times 1.0 + 16.0 = \qty{46.0}{g/mol}$، إذن
$1000 / 46.0 = \qty{21.7}{mol}$. فللماء الكمية الأكبر، بعامل
$55.6 / 21.7 \approx 2.6$، وهو ببساطة $46.0 / 18.0$.
\end{solution}

\begin{solution}{exo:g10:the-mole:12}
الذهب: $\qty{3.75}{g}$، $n = 3.75 / 197.0 = \qty{1.90e-2}{mol}$، أي
\num{1.15e22} \omterm{def:g7:atoms-and-molecules:atom}{ذرة}. النحاس: \qty{1.25}{g}،
$n = 1.25 / 63.5 = \qty{1.97e-2}{mol}$، أي \num{1.19e22} \omterm{def:g7:atoms-and-molecules:atom}{ذرة}.
والمفاجئ أن الخاتم يحوي من \omterm{def:g7:atoms-and-molecules:atom}{ذرات} النحاس أكثر قليلًا من \omterm{def:g7:atoms-and-molecules:atom}{ذرات} الذهب:
\omterm{def:g7:atoms-and-molecules:atom}{فذرة} الذهب أثقل من \omterm{def:g7:atoms-and-molecules:atom}{ذرة} النحاس بنحو ثلاث مرات.
\end{solution}

\begin{solution}{exo:g10:the-mole:13}
\begin{enumerate}
  \item $V = 8.0 \times 6.0 \times 3.0 = \qty{144}{m^3} =
        \qty{1.44e5}{L}$؛ $n = \num{1.44e5} / 24.1 = \qty{5.98e3}{mol}$؛
        $N = \num{5.98e3} \times \num{6.02e23} = \num{3.6e27}$ \omterm{def:g7:atoms-and-molecules:molecule}{جزيء}.
  \item $M = (78 \times 28.0 + 21 \times 32.0 + 1 \times 40.0) / 100 =
        2896 / 100 = \qty{29.0}{g/mol}$.
  \item $m = \num{5.98e3} \times 29.0 = \qty{1.73e5}{g}$، أي نحو
        \qty{173}{kg}: \omterm{def:g4:air-a-mixture-of-gases:air}{فهواء} القاعة يزن ما يزنه شخصان أو ثلاثة
        بالغون.
\end{enumerate}
\end{solution}

\begin{solution}{exo:g10:the-mole:14}
\begin{enumerate}
  \item $n = 1000 / 28.1 = \qty{35.6}{mol}$، إذن
        $N = 35.6 \times \num{6.02e23} = \num{2.14e25}$ \omterm{def:g7:atoms-and-molecules:atom}{ذرة}.
  \item $m = \qty{1.000}{kg} / \num{2.14e25} = \qty{4.67e-26}{kg}$.
  \item كتلة الكرة وكتلتها المولية تعطيان كمية مادتها،
        $n = m / M$. فإذا عُدّت \omterm{def:g7:atoms-and-molecules:atom}{الذرات} $N$ عدًّا مستقلًا (من حجم الكرة
        والمسافة بين \omterm{def:g7:atoms-and-molecules:atom}{الذرات} في البلورة)، كانت النسبة $N / n$ هي \omterm{def:g10:the-mole:avogadro}{ثابت
        أفوغادرو}.
\end{enumerate}
\end{solution}

\begin{solution}{exo:g10:the-mole:15}
$M = 0.758 \times 35.0 + 0.242 \times 37.0 = 26.53 + 8.95 =
\qty{35.5}{g/mol}$: قيمة الجدول.
\end{solution}

\begin{solution}{pb:g10:the-mole:1}
\textbf{1.} $M(\ce{H2O}) = 2 \times 1.0 + 16.0 = \qty{18.0}{g/mol}$.

\textbf{2.} $m = 250 \times \qty{1.0}{g} = \qty{250}{g}$.

\textbf{3.} $n = 250 / 18.0 = \qty{13.9}{mol}$.

\textbf{4.} في كل \omterm{def:g7:atoms-and-molecules:molecule}{جزيء} ذرتا هيدروجين \omterm{def:g7:atoms-and-molecules:atom}{وذرة} أكسجين واحدة:
\qty{27.8}{mol} من \omterm{def:g7:atoms-and-molecules:atom}{ذرات} الهيدروجين، و\qty{13.9}{mol} من \omterm{def:g7:atoms-and-molecules:atom}{ذرات} الأكسجين.

\textbf{5.} $N = 13.9 \times \num{6.02e23} = \num{8.36e24}$ \omterm{def:g7:atoms-and-molecules:molecule}{جزيء}.

\textbf{6.} $m = \qty{18.0}{g/mol} / \qty{6.02e23}{mol^{-1}} =
\qty{2.99e-23}{g}$.

\textbf{7.} $\num{8.36e24} \times \qty{2.99e-23}{g} = \qty{250}{g}$، وهي
كتلة الماء في الكأس، كما ينبغي.

\textbf{8.} $\num{8.36e24} / \num{3.16e7} = \num{2.6e17}$ سنة: عدد سنوات
بسبعة عشر صفرًا، أبعد بكثير من أي عدّ ممكن.

\textbf{9.} $\qty{1.335e9}{km^3} \times \num{e9} = \qty{1.335e18}{m^3}$،
ومع $\qty{1}{m^3} = \qty{1000}{L}$ نحصل على \qty{1.335e21}{L}.

\textbf{10.} $\num{8.36e24} / \num{1.335e21} = \num{6.26e3}$ \omterm{def:g7:atoms-and-molecules:molecule}{جزيء} معلَّم
في كل لتر.

\textbf{11.} $\num{6.26e3} \times 0.250 = \num{1.57e3}$ \omterm{def:g7:atoms-and-molecules:molecule}{جزيء} معلَّم
في الكأس الثانية.

\textbf{12.} الكتلة $\num{1.335e21} \times \qty{1.0}{kg} =
\qty{1.335e21}{kg} = \qty{1.335e24}{g}$؛ و$n = \num{1.335e24} / 18.0 =
\qty{7.42e22}{mol}$.

\textbf{13.} $13.9 / \num{7.42e22} = \num{1.87e-22}$.

\textbf{14.} $\num{1.87e-22} \times \num{8.36e24} = \num{1.57e3}$: الجواب
نفسه الذي في السؤال 11، كما يجب.

\textbf{15.} $\num{1.335e21} / 0.250 = \num{5.34e21}$ كأس.

\textbf{16.} $\num{8.36e24} / \num{5.34e21} \approx 1570$. فالكأس تحوي من
\omterm{def:g7:atoms-and-molecules:molecule}{الجزيئات} نحو 1600 مرة أكثر مما تحوي المحيطات من الكؤوس؛ وإذا انتشرت \omterm{def:g7:atoms-and-molecules:molecule}{جزيئات}
كأس واحدة بانتظام بقي منها نحو 1600 في كل كأس من ماء
المحيط.

\textbf{17.} $1 / \num{6.26e3} = \qty{1.6e-4}{L} = \qty{0.16}{mL}$، أي نحو
ثلاث قطرات.

\textbf{18.} حجم المحيطات تقدير معروف بثلاثة أرقام أو أربعة في أحسن الأحوال،
وقد قُرّب؛ وماء البحر ليس ماءً نقيًا (فهو يحوي أملاحًا)؛ والكأس ليست
\qty{250}{mL} بالضبط؛ والاختلاط في المحيطات كلها لن يكون أبدًا منتظمًا
تمامًا. فلا يُوثق إلا برتبة القدر.

\textbf{19.} يعود نحو 1600 \omterm{def:g7:atoms-and-molecules:molecule}{جزيء} من الكأس الأولى في
الثانية.
\end{solution}
